\begin{table}[t]
\centering
\small
\caption{Interpretation of close strong-ablation comparisons. The table clarifies what each comparison supports and what it does not support.}
\label{tab:strong-ablation-interpretation}
\begin{tabularx}{\linewidth}{p{0.27\linewidth}p{0.23\linewidth}X}
\toprule
Comparison & Observed outcome & Supported interpretation \\
\midrule
Full vs. flat knowledge LTR, no graph & No significant difference in MRR@10, Recall@10, or nDCG@10 & Flat retrieval, semantic, and biomedical knowledge features explain most of the aggregate gain; the full model should not be claimed to dominate this ablation. \\
Full vs. Pairwise graph LTR & No significant difference at $k=10$; nDCG@10 is a near-significant positive trend & The hypergraph is a useful structural variant, but current top-10 evidence does not prove a clear aggregate advantage over pairwise graph decomposition. \\
Full vs. Remove MeSH hierarchy & No significant difference; the ablation has slightly higher Recall@10 and nDCG@10 & MeSH hierarchy features are best presented as interpretable medical constraints, not as an isolated source of statistically proven aggregate gain. \\
Full vs. Remove hypergraph diffusion/centrality & No significant difference; the ablation has slightly higher MRR@10 and nDCG@10 & Diffusion and centrality features aid transparency and case-level explanation, but the current test does not prove that they independently increase overall ranking metrics. \\
\bottomrule
\end{tabularx}
\end{table}
